\section{主要符号及约定}\label{app:notation}

\begin{longtable}{p{31mm}p{112mm}}
\toprule
符号 & 含义\\
\midrule
\endhead
$m$ & 训练样本数量。\\
$p$ & 参数维数；在具体任务中应与特征维数的约定相区分。\\
$(x_i,y_i)$ & 第 $i$ 个输入与目标观测。\\
$D$ & 训练数据集。\\
$\mathcal F$ & 候选预测函数或模型族。\\
$f_\theta$ & 由参数 $\theta$ 指定的预测函数。\\
$L(y,\hat y)$ & 单个样本的损失函数。\\
$\widehat R_D(f)$ & 数据集 $D$ 上的经验风险。\\
$R(f)$ & 在目标分布下的期望风险。\\
$J(\theta)$ & 需要最小化的目标函数；其尺度按上下文给出。\\
$X\in\R^{m\times p}$ & 设计矩阵。与截距相关的常数列可包含在内。\\
$\theta,\beta$ & 模型参数向量，选用何者应在相应公式附近说明。\\
$\nabla J,\nabla^2J$ & 目标函数的梯度与 Hessian 矩阵。\\
$\eta_k$ & 第 $k$ 次迭代所采用的学习率。\\
$\pi(x)$ & 逻辑回归中正类的条件概率。\\
$TP,FP,FN,TN$ & 二分类混淆矩阵的四种计数。\\
$\E,\Var,\Bias$ & 期望、方差与偏差。\\
$\norm{v}_2$ & 向量 $v$ 的欧氏范数。\\
\bottomrule
\end{longtable}

除非另行说明，文中向量按列向量处理，转置记作 $\trans$。对符号相同但假设不同的结论，应在首次使用时重新说明适用范围。文中所称“最优解”通常是训练目标意义上的最优解，而不是泛化风险意义上的最优解。
